\section{Discussion}
\label{sec:discussion}

\subsection{What the framework says to build}

Read as engineering advice, the decomposition gives a short checklist that can
be run in an afternoon on \nCalib{} labelled tasks, before any orchestration
code is written.

\begin{enumerate}[leftmargin=1.4em,itemsep=2pt,topsep=3pt]
\item \textbf{Measure the headroom.} Sample $k$ answers per task and compute
  $\cov(k)-\pone$. If it is small, no architecture will help, and the effort
  belongs in a better generator or a better prompt rather than in
  orchestration.
\item \textbf{Measure the gap $\gap$.} Show the model one correct and one
  incorrect candidate and ask which is right. If $\gap \approx 0$, a
  language-model critic will not convert the headroom, no matter how the
  critique is prompted or how many rounds it runs.
\item \textbf{Look for an external check first.} A program that validates the
  task's own stated constraints is the highest-efficiency selector available
  and, unlike a judge, consumes no model tokens. Much of the benefit attributed
  to multi-agent design in practice is really the benefit of having a checker.
\item \textbf{Price the alternative.} Compare the predicted gain against what
  one agent achieves with the same budget spent serially. This is the
  comparison that decides whether the architecture is worth its complexity.
\end{enumerate}

The uncomfortable implication for multi-agent design is that the most common
pattern---adding a language-model critic or reviewer agent---is the one our
framework is least optimistic about, because it depends entirely on $\gap$
being large for the specific model and task at hand, and $\gap$ is exactly what
practitioners do not currently measure.

\subsection{Reconciling conflicting reports}

The framework explains why the literature disagrees without anyone being wrong.
Studies that report large multi-agent gains tend to use tasks with substantial
coverage headroom and either a strong external verifier or an answer
distribution peaked on the truth; studies that report little benefit tend to use
tasks where verification is as hard as generation, or compare against stronger
baselines. Because $\gap$ and $\plur$ vary across both tasks \emph{and} models,
two papers can run the same architecture on different benchmarks and reach
opposite conclusions that are both correct. The prescription is not to settle
the argument in general but to measure the two diagnostics in each new setting.

\subsection{Limitations}
\label{sec:limitations}

\paragraph{Model scale.} Our experiments use open models at 0.5B and 1.5B parameters, run locally. This was a deliberate choice---it makes the study
exactly reproducible at zero marginal cost, which matters for a paper whose
central claim is about compute accounting---but it is also the main limit on
what we can claim. Frontier models are better verifiers, so $\gap$ should be
larger for them and selection efficiency correspondingly higher. Our scale analysis (Section~\ref{sec:results}) contrasts the two scales we have, which shows the diagnostics are model-dependent rather than task-intrinsic but establishes no trend: two points determine a line and test nothing. Extrapolating to models two orders of magnitude larger is not something our data licenses, and we make no such claim. What we do claim is
scale-independent: the \emph{decomposition} is an identity, the need for
compute-matched baselines is a matter of accounting, and the diagnostics remain
measurable at any scale. Whether the specific numbers transfer is an empirical
question we cannot answer here.

\paragraph{The latent-score model is a simplification.} Definition~\ref{def:latent}
assumes verifier scores are conditionally independent across candidates and
equally dispersed for correct and incorrect ones. Real judges violate both:
they see all candidates at once, and a confident wrong answer may attract a
higher score than a hedged right one. We report where the prediction holds and
where it breaks rather than claiming the model is correct.

\paragraph{Task families are a sample, not a census.} Six families cannot span
every kind of verification asymmetry. Our synthetic families are exactly
checkable by construction, which is what makes them useful anchors, but it also
makes them unrepresentative of open-ended work where no ground truth exists at
all---precisely the regime where much multi-agent deployment happens, and where
our framework predicts the least benefit but can offer no measurement.

\paragraph{Architectures are a sample too.} We test the designs that dominate
the literature---voting, judging, debate, refinement---but not tool-using
agents, long-horizon planning, or role-specialised pipelines where agents solve
genuinely different sub-problems. Our framework applies to the selection step of
any such system; it does not model the decomposition step.

\paragraph{Single-answer tasks.} The decomposition assumes a task has a
selectable answer. It does not directly cover systems whose agents each
contribute a different part of one artifact, which is how many deployed
pipelines---including the one in our case study---actually work. We address that
gap only structurally, by classifying which stages admit an external check.

\section{Conclusion}

We set out to answer a question practitioners face constantly and currently
answer by guesswork: is a multi-agent architecture worth building for this task?
The answer we arrive at is that the question is well-posed and cheaply
decidable. An ensemble's accuracy is its coverage times its selection
efficiency, selection efficiency is predictable from a pairwise probe that costs
a few dozen model calls, and the honest baseline is one agent given the same
tokens. Measured this way, multi-agent structure stops looking like a general
technique and starts looking like what it is: an effective way to exploit an
asymmetry between generating and checking, and an expensive way to spend tokens
when that asymmetry is absent.
